\documentclass[12pt,a4paper]{article}
\usepackage[utf8]{inputenc}
\usepackage{geometry}
\geometry{a4paper, margin=1in}

\title{\textbf{How We Calculate the Buyer Intent Score \\ \large A Simple Guide for Investors}}
\author{Thhiya Data Science}
\date{\today}

\begin{document}

\maketitle

\section*{The Big Picture}
Our platform uses an advanced AI "brain" (Google Gemini) to read millions of data points across the internet. It acts like a highly intelligent, speed-reading researcher. When we want to evaluate a company, the AI automatically reads its public reviews, the latest news, and overall online reputation to calculate a dynamic "Buyer Intent Score." 

Think of this score as a measure of how "hot" or "in-demand" a company is right now.

\section*{The Core Formula (What the AI Looks At)}
The AI does not guess the score. We strictly instruct it to grade the company based on four specific pillars of health. We can think of the final score as a weighted average of these four areas:

\textbf{Final Intent Score = (Market Momentum) + (User Sentiment) + (Innovation) + (Transparency)}

\begin{itemize}
    \item \textbf{Market Momentum:} Is the company growing rapidly? Are people talking about it more than usual?
    \item \textbf{User Sentiment:} Are the reviews online overwhelmingly positive? Are customers happy?
    \item \textbf{Feature Innovation:} Is the company releasing new, exciting tech (like AI integration), or are they falling behind competitors?
    \item \textbf{Transparency:} Are their prices and services clear and honest, or are customers frustrated by hidden fees?
\end{itemize}

Based on how strong a company is in each of these four areas, the AI calculates a final number.

\section*{The "Guardrails" (Ensuring Accurate Data)}
To make sure our software never breaks and the investors get reliable data, our engineering system places strict "guardrails" on top of the AI's calculation. 

Mathematically, our system enforces this logic:
\[ \textbf{Final Score} = \textbf{Ensure it is between 0 and 10} \]

\begin{itemize}
    \item \textbf{The Cap:} The score can never be higher than 10 or lower than 0.
    \item \textbf{The Safety Net:} If the AI fails to generate a score (e.g., if a company is so new that there is no data on the internet), our system automatically assigns a neutral, middle-ground score of \textbf{5}. This ensures our app never crashes and always displays a realistic baseline.
\end{itemize}

\section*{The Proof (Why Should We Trust It?)}
We don't just output a number. To prove the score is accurate, the AI is legally required by our system to fetch evidence. Alongside every score, our system extracts and saves:
\begin{enumerate}
    \item Exactly \textbf{5 distinct positive reviews} summarizing what people love.
    \item Exactly \textbf{5 distinct negative reviews} summarizing the genuine complaints.
    \item A 3-to-5 word justification for \textit{each} of the four pillars (e.g., "Rapidly growing user base").
\end{enumerate}
This means every single 0-10 score presented to users or investors is instantly backed up by readable, real-world proof.

\end{document}
